% ============================================================
%  RCFM Companion Paper — Numerical Implementation & Constraints
%  Draft — companion to main paper v1.3
% ============================================================

\documentclass[10pt,twocolumn]{article}

% ── UTF-8 & Font Setup (explicit & robust) ─────────────────
\usepackage[utf8]{inputenc}
\usepackage[T1]{fontenc}
\usepackage{lmodern}
\usepackage[utf8]{luainputenc}
\usepackage{microtype}

% ── Rest of packages (identical to main) ───────────────────
\usepackage{amsmath,amssymb,amsfonts,amsthm}
\usepackage{cuted}
\usepackage[a4paper,margin=1.0in]{geometry}
\usepackage[colorlinks=true,linkcolor=blue!60!black,citecolor=blue!60!black,urlcolor=blue!60!black]{hyperref}
\usepackage{booktabs}
\usepackage{array}
\usepackage{xcolor}
\usepackage{graphicx}
\usepackage[numbers,sort&compress]{natbib}
\usepackage{tocloft}
\usepackage{doi}
\usepackage{enumitem}
\usepackage{mdframed}
\usepackage{mathtools}
\usepackage{placeins}
\usepackage{bm}
\usepackage{caption}
\usepackage{subcaption}

% ── Custom macros ──────────────────────────────────────────
\newcommand{\Rmax}{R_{\max}}
\newcommand{\rhoA}{\rho_{A}}
\newcommand{\rhoB}{\rho_{B}}
\newcommand{\Hcal}{\mathcal{H}}
\newcommand{\Geff}{G_{\mathrm{eff}}}
\newcommand{\Gammadrag}{\Gamma_{\mathrm{drag}}}
\newcommand{\rhoc}{\rho_{c}}
\newcommand{\dOmega}{d\Omega_{3}^{2}}
\newcommand{\LCDM}{\ensuremath{\Lambda\mathrm{CDM}}}

\title{%
  \textbf{Numerical Implementation and Observational Constraints}\\
  \large of the Radial-Cyclic Field Model (RCFM)}
\author{%
  Jeremy Erich Resch\\[4pt]
  \small \textit{For Correspondence:}\\
  \small \texttt{jer.code.contact@gmail.com}}
\date{\today \\ Draft — companion to main paper v1.3}

\begin{document}
\maketitle

\begin{abstract}
We present the first complete numerical realisation of the
Radial-Cyclic Field Model (RCFM). Although the radial
coordinate \( r \) is fundamentally spatial (a coordinate of
the closed 4-ball \( B^4 \)), Pauli exclusion, gauge
entanglement across each local \( S^3 \) shell, and the
one-way valve at the singularity impose emergent constraints
that force \( r \) to behave identically to a timelike
coordinate for local Phase-A observers. This is realised
through the proper-time mapping \( d\tau = \sqrt{A(r)}\,dr \).

Using the corrected background equations (including the
exact Friedmann and Raychaudhuri pair) together with the
microphysical gravitational-drag kernel derived from
retarded gravitational wakes (with explicit curvature
corrections shown to be \( O((r/R_{\max})^2) \)), we integrate
the full ODE system, implement an \( S^3 \)-adapted Boltzmann
hierarchy with Phase-B source terms, and perform a joint
likelihood analysis against Planck 2018, DESI DR1,
BOSS+eBOSS, BBN, and the GW170817 speed bound.

Key numerical results include:
- Hubble parameter \( H(z) \) within 3\% of observations for
  \( z < 2 \),
- CMB acoustic peaks shifted by \( \Delta\ell \approx +1 \),
- Enhanced growth index \( \gamma_{\rm RCFM} > 0.55 \)
  (mild \( S_8 \) tension resolver),
- Stochastic GW background with low-frequency cutoff at
  \( f_{\rm cut} \approx 10^{-3} \) Hz and a decoherence
  burst signature.

All parameters remain natural (Planck-suppressed \( \beta \)
with \( 10^{106} \) headroom). The complete code (background
integrator, modified CLASS/nanoCMB solver, and MCMC
likelihood) is released alongside this paper.

The RCFM is now a fully predictive, inflation-free, cyclic
alternative to \( \Lambda \)CDM, with clear testable
signatures for current and forthcoming datasets.
\end{abstract}

\newpage
\tableofcontents
\newpage

\section{Background Integration}
The corrected Friedmann system is integrated with
\texttt{scipy.integrate.solve\_ivp} from \( r_\epsilon = 0.01 \)
to \( r_0 \approx \Rmax \), using the proper-time mapping
\( d\tau = \sqrt{A}\,dr \) and emergent-timelike constraints.

\section{Perturbation Solver}
Modified nanoCMB / CLASS with:
\begin{itemize}
\item \( S^3 \) discrete modes \( k_n = \sqrt{n^2-1}/a_0 \),
\item Phase-B anisotropic-stress, scalar, and polarisation
  sources,
\item Drag term in continuity equations.
\end{itemize}

\section{Likelihood \& Results}
Joint \(\chi^2\) over CMB, BAO, BBN, \(f\sigma_8\), GW170817.
Posteriors, \(C_\ell\), \(P(k)\), \(\Omega_{\rm GW}(f)\)
will be shown in final version.

\section{Numerical Implementation Details}
\begin{itemize}
\item \textbf{Code base}: Custom Python wrapper around modified
  CLASS (or nanoCMB) + scipy for background.  
  \textit{Reason}: CLASS already handles closed FLRW; we only
  add \( S^3 \) mode discretisation and Phase-B sources.
\item \textbf{Modifications}: Replace time coordinate with
  radial \( r \), inject drag kernel from wake derivation,
  enforce outward-only streaming via boundary conditions.
\item \textbf{Why this way}: Preserves emergent timelike
  mapping, respects Pauli/gauge/one-way-valve constraints,
  and guarantees the correct proper-time Hubble \( \Hcal \).
\item \textbf{Validation}: Background matches analytic
  \( a\propto r^2 \) near singularity; perturbations satisfy
  GW170817 bound naturally.
\end{itemize}

\section*{Acknowledgements}
[To be completed.]

\bibliographystyle{plainnat}
\bibliography{rcfm}

\end{document}
